\section{Experimental Methodology}

The evaluation should map each headline claim to at least one measurement and
one competing explanation.  Table~\ref{tab:questions} provides a compact plan.
End-to-end measurements establish practical relevance; ablations attribute the
effect; sensitivity analyses delimit the operating region; and adversarial or
failure workloads identify where the design should not be used.

\begin{table}[H]
  \caption{Claim-to-evidence plan.}
  \label{tab:questions}
  \centering
  \small
  \setlength{\tabcolsep}{3pt}
  \begin{tabular}{L{0.20\columnwidth}L{0.29\columnwidth}L{0.37\columnwidth}}
    \toprule
    Question & Measurement & Reporting rule \\
    \midrule
    Efficiency & Throughput, median and tail latency & Same hardware, semantics,
    load, and tuning budget \\
    Attribution & Component ablations and profiles & Report absolute values as
    well as relative change \\
    Robustness & Data scale, skew, concurrency, and failures & Show the full
    range, including regressions \\
    Cost & CPU, memory, storage, network, and dollars & State accounting window
    and pricing date \\
    \bottomrule
  \end{tabular}
\end{table}

\subsection{Baselines and configuration}

Baselines should represent the strongest relevant published and production
approaches, not merely implementations that are convenient to run.  Apply a
documented and comparable tuning budget to every system.  Record source
revision, compiler, flags, dependencies, machine type, storage medium, network
topology, and operating-system settings.  When a proprietary baseline prevents
inspection, clearly separate observed behavior from inferred causes.

\subsection{Workloads and metrics}

Use standard benchmarks where they capture the claimed setting and supplement
them with targeted workloads where they do not.  Describe data generation,
scale, skew, arrival process, warm-up, cache state, concurrency, and termination
conditions.  Report distributions or meaningful quantiles in addition to
means.  For repeated trials $x_1,\ldots,x_k$, state $k$, the aggregation
procedure, and uncertainty; for example, a 95\% confidence interval around a
mean should be accompanied by the resampling or parametric method used to
compute it.

\subsection{Analysis protocol}

Fix the analysis protocol before inspecting the final results.  Distinguish
confirmatory questions from exploratory observations, retain failed runs with
documented reasons, and avoid selecting only favorable workload slices.  A
machine-readable result table should contain raw measurements and identifiers
that join each row to its configuration and execution log.
